% =====================================================================================================
\chapter{Machine-learning methods}\label{ch:methods}
% =====================================================================================================

Every model used in the project is defined here once, in the form in which it was actually run: scikit-learn 1.8.0
(\file{requirements-lock.txt}) with the settings of \file{scripts/config.py} and the wrappers \code{V4Model}
(\file{v4lib.py}), \code{C2Model} (\file{v7lib.py}) and the pipelines of \file{07\_train\_eval.py} and \file{v3lib.py}.
Equations describe the scikit-learn implementation, not a generic textbook variant. \Cref{tab:modelinventory} is the
inventory; \cref{app:hyper} collects all hyper-parameters.

\begin{table}[htbp]
\centering\small
\caption[Inventory of models]{Inventory of the models and where they are used. Role: P = model of a pre-registered
primary comparison; B = baseline; D/E = descriptive/exploratory only.}
\label{tab:modelinventory}
\begin{tabularx}{\textwidth}{@{}p{3.2cm}p{3.6cm}Xc@{}}
\toprule
Model & Versions & Code & Role \\
\midrule
Dummy (prior) & v1, v2 & \file{07\_train\_eval.py} & B \\
Logistic Regression & v1--v13 & \file{07}, \code{v3lib.MODELS}, \code{V4Model}, \code{C2Model}, \code{CCTLR.lr} & P (v2--v11), B \\
Random Forest & v1--v13 & as above & P (v8d, v12b, v13a), D \\
ExtraTrees & v4a & \code{V4Model} & E \\
HistGradientBoosting & v4a & \code{V4Model} & E \\
RBF SVM & v4a, v7c2; v7c1 (Platt, 3-class) & \code{V4Model}, \code{C2Model}, \file{37\_v7c1\_repro.py} & E / reproduction \\
kNN classifier & v4a, v7c2, v8d; v7c1 (3-class) & \code{V4Model}, \code{C2Model}, \file{37} & E / reproduction \\
QDA & v4a & \code{V4Model} & E \\
MLP (seed ensemble) & v4a & \code{V4Model} & E \\
Nested auto-selection & v4a 1d/1e & \code{v4lib.nested\_fold} (\cref{sec:nested}) & P (v4) \\
CCTLR & v6c & \code{v6lib.CCTLR} & D (new method) \\
kNN regression & v6d, v9b--v13b & \code{v6lib.source\_residuals} & inside the tests \\
Frozen $\nu^\ast$ rule & v11, v13 & \file{57\_v11\_analysis.py}, \file{63\_v13\_analysis.py} & D \\
\bottomrule
\end{tabularx}
\end{table}

% -----------------------------------------------------------------------------------------------------
\section{Shared components}\label{sec:shared}
% -----------------------------------------------------------------------------------------------------

\subsection{Standardisation}\label{sec:scaler}
For the scale-sensitive models (LR, SVM, kNN, QDA, MLP, the colour part of CCTLR and the kNN regression) each feature
is standardised with the training rows $\mathcal T$ of the current fit only (\code{StandardScaler}):
\begin{equation}
\mu_\ell = \frac{1}{|\mathcal T|}\sum_{i\in\mathcal T}x_{i\ell},\qquad
\sigma_\ell = \Bigl(\frac{1}{|\mathcal T|}\sum_{i\in\mathcal T}(x_{i\ell}-\mu_\ell)^2\Bigr)^{1/2},\qquad
z_{i\ell} = \frac{x_{i\ell}-\mu_\ell}{\sigma_\ell}\quad(\sigma_\ell=0\Rightarrow\text{scale }1),\label{eq:scaler}
\end{equation}
and the same $(\mu_\ell,\sigma_\ell)$ are applied to test rows. Tree models (RF, ExtraTrees, HistGB) use raw features.
Standardisation is learned; the per-spectrum normalisations of \cref{sec:reps} (B, H, \dots) are not.

\subsection{Class-balanced weights}\label{sec:weights}
The BH:NS ratio is 7:24 sources in the main sample. With $n$ training rows of which $n_c$ are in class $c$,
\code{class\_weight='balanced'} gives every row of class $c$ the weight
\begin{equation}
w^{\mathrm{bal}}_i = \frac{n}{2\,n_{y_i}},\qquad \sum_{i:y_i=c}w^{\mathrm{bal}}_i = \frac n2 .\label{eq:balanced}
\end{equation}
In scikit-learn these class weights multiply any sample weights inside the loss (LR, SVM, HistGB) or inside the impurity
(trees). For QDA (no weights) the same balance is imposed through equal priors, for kNN through a prior correction of the
output (\cref{eq:knnprior}), and for the MLP through $w^{\mathrm{bal}}$ passed as sample weights.

\paragraph{Source-equal weights (v4a 1c, v4b1, v5 S3, v6 S3).} When sources contribute very different numbers of
observations (v4b1: 6 to 1329 per source), balanced weights let the most-observed sources dominate. With $S_c$ the number
of training sources in class $c$ and $n_k$ the training rows of source $k$ (\code{v4lib.source\_equal\_weights}):
\begin{equation}
w^{\mathrm{src}}_i = \frac{n}{2\,S_{y_i}\,n_{g_i}},\qquad \sum_{i:g_i=k} w^{\mathrm{src}}_i = \frac{n}{2S_{c}}\ \ (k\in\text{class }c),\qquad
\sum_{i:y_i=c}w^{\mathrm{src}}_i=\frac n2 .\label{eq:sourceequal}
\end{equation}
Every source of a class has the same total weight, both classes sum to $n/2$ (the same overall scale as balanced weights),
and $w^{\mathrm{src}}=w^{\mathrm{bal}}$ exactly when all sources have the same $n_k$ (unit test in \file{test\_v4lib.py}).
When $\wv^{\mathrm{src}}$ is passed, \code{class\_weight} is set to \code{None}.

\begin{cautionbox}{Implementation fix: scale of the source-equal weights (2026-10-01, after the v4a results)}
The first version normalised each class to total weight 1 (total 2 per fit) instead of $n/2$. Because the LR objective
multiplies the summed weighted loss by $C$ (\cref{eq:lrobj}), this acted like an $\approx n/2$ times stronger L2 penalty
(scores squeezed into 0.4--0.6); the HistGB internal \code{min\_hessian\_to\_split} was also affected (RF is scale-invariant).
The fix rescaled to $n/2$ per class; only v4a parts 1a--1c were re-run: 1a and 1b were bit-identical, 12 combinations of 1c
changed, and no conclusion changed (pre-fix results in commit \texttt{9120afb}). v4b1 was run only with the fixed weights.
\end{cautionbox}

\subsection{Score conventions}
Every model returns a BH score $s(\xv)\in[0,1]$: \code{predict\_proba}$[:,1]$ for LR, RF, ExtraTrees, HistGB, QDA;
the prior-corrected vote fraction for kNN; the seed-averaged probability for the MLP; the logistic function of the
decision value for the SVM and CCTLR. Scores are rankings and thresholdable quantities, not calibrated probabilities,
except where proper scoring rules are explicitly evaluated (v6g).

% -----------------------------------------------------------------------------------------------------
\section{Dummy classifier}\label{sec:dummy}
% -----------------------------------------------------------------------------------------------------
\textbf{Role.} Baseline in v1--v2 (\file{07\_train\_eval.py}). \textbf{Definition.} \code{DummyClassifier(strategy='prior')}
returns the training class frequencies, $s(\xv)=n_1/n$ for every $\xv$, and predicts the majority class. \textbf{Why it
matters.} Under LOSO the majority class of the training fold depends on which source is held out: with 4 BH and 4 NS
sources of 15 observations each (v1), leaving out a BH leaves 45 BH : 60 NS, so the held-out BH is predicted NS, and vice
versa; the Dummy balanced accuracy is therefore 0, not 0.5. With 7 BH and 24 NS sources (v2) NS is always the majority
and the Dummy balanced accuracy is 0.5. ``Chance level'' under LOSO is thus not a fixed number, one reason why AUC is
also reported.

\begin{algorithm}[htbp]
\caption{Dummy (prior) classifier}\label{alg:dummy}
\begin{algorithmic}[1]
\Require training labels $\yv_\mathcal T$; test rows $\mathcal V$
\State $\pi\gets\frac{1}{|\mathcal T|}\sum_{i\in\mathcal T}y_i$ \Comment{no features; only hyper-parameter \code{strategy='prior'}}
\State \Return $s_i=\pi$ and $\hat y_i=\indic{\pi\ge0.5}$ for every $i\in\mathcal V$
\end{algorithmic}
\end{algorithm}


% -----------------------------------------------------------------------------------------------------
\section{Logistic regression}\label{sec:lr}
% -----------------------------------------------------------------------------------------------------
\textbf{Role.} The workhorse: H-LR (two colours + LR) is the v2 reference against which v4 tests every alternative;
H+T-LR is the primary model of v5, v6a, v8a/b, v9a, v10a, v11a. \textbf{Model.} On standardised inputs $\vect z$,
\begin{equation}
s(\xv) = \sigma(\vect w^\top\vect z + b),\qquad \sigma(u)=\frac{1}{1+e^{-u}} .\label{eq:lrmodel}
\end{equation}
\textbf{Objective} (\code{LogisticRegression(C=1, penalty='l2', solver='lbfgs', max\_iter=5000)}):
\begin{equation}
(\hat{\vect w},\hat b) = \argmin_{\vect w,b}\ \frac12\|\vect w\|_2^2 + C\sum_{i\in\mathcal T} \omega_i\Bigl[-y_i\log s(\xv_i)-(1-y_i)\log\bigl(1-s(\xv_i)\bigr)\Bigr],\label{eq:lrobj}
\end{equation}
with $\omega_i = w^{\mathrm{bal}}_i$ (\cref{eq:balanced}) or $w^{\mathrm{src}}_i$ (\cref{eq:sourceequal}); the intercept is
not penalised. scikit-learn minimises the equivalent $\frac{1}{\Omega}\sum_i\omega_i\ell_i + \frac{1}{2C\Omega}\|\vect w\|^2$
($\Omega=\sum_i\omega_i$) with L-BFGS. $C=1$ means the penalty weight $1/2$ competes with a weighted log-loss summed over
$\approx n$ rows: for $n\approx400$ and 2--6 features the penalty is weak. \textbf{Decision boundary.} In raw colours the
0.5 boundary of H-LR is the line $\sum_\ell (w_\ell/\sigma_\ell)h_\ell + b - \sum_\ell w_\ell\mu_\ell/\sigma_\ell = 0$
(used to draw \cref{fig:colourcolour}). \textbf{Assumptions.} Log-odds linear in the standardised features; a single
monotone effect per feature. This is the main reason LR can use timing only as ``more variability $\Rightarrow$ more BH''
everywhere in the colour plane (motivating CCTLR, \cref{sec:cctlr}), and it fails on MAXI where BHs lie on both sides of
the NS colour track (\cref{ch:phase4}).

\begin{algorithm}[htbp]
\caption{LR fit and score (one fold)}\label{alg:lr}
\begin{algorithmic}[1]
\Require $\Xm_\mathcal{T},\yv_\mathcal{T}$, optional $\wv$; test rows $\Xm_\mathcal{V}$
\State $(\vect\mu,\vect\sigma)\gets$ \cref{eq:scaler} on $\Xm_\mathcal{T}$; $\vect Z_\mathcal{T},\vect Z_\mathcal{V}$ standardised
\State $\omega_i\gets w^{\mathrm{bal}}_i$ if $\wv$ is None else $w_i$
\State $(\vect w,b)\gets$ L-BFGS minimisation of \cref{eq:lrobj} (max.\ 5000 iterations)
\State \Return $s_i = \sigma(\vect w^\top\vect z_i+b)$ for $i\in\mathcal V$
\end{algorithmic}
\end{algorithm}

\begin{table}[htbp]
\centering\small
\caption[Logistic-regression hyper-parameters]{Logistic-regression hyper-parameters (\code{LR\_PARAMS} in \file{config.py}).}\label{tab:hp-lr}
\begin{tabularx}{\textwidth}{@{}llp{3.0cm}X@{}}
\toprule
Parameter & Value & Fixed / tuned & Grid (where tuned) \\
\midrule
\code{C} & 1.0 & fixed (all versions) & v4a 1d/1e: $\{0.01, 0.1, 1, 10\}$ (\code{V4\_GRID['LogReg']}) \\
\code{penalty}, \code{solver} & L2, lbfgs & fixed & -- \\
\code{class\_weight} & balanced (or \code{None} with $\wv^\mathrm{src}$) & fixed & -- \\
\code{max\_iter} & 5000 & fixed & -- \\
\bottomrule
\end{tabularx}
\end{table}

% -----------------------------------------------------------------------------------------------------
\section{Random forest}\label{sec:rf}
% -----------------------------------------------------------------------------------------------------
\textbf{Role.} Second reference model in every version; primary model of v8d (MAXI), v12b (chosen post hoc after
v9--v11) and v13a (chosen before the v13 data). \textbf{Model.} An average of $B$ classification trees,
\begin{equation}
s(\xv) = \frac1B\sum_{b=1}^{B}\hat p_b(\xv),\qquad
\hat p_b(\xv)=\frac{\sum_{i\in\mathrm{leaf}_b(\xv)}\omega^{(b)}_i y_i}{\sum_{i\in\mathrm{leaf}_b(\xv)}\omega^{(b)}_i},\label{eq:rf}
\end{equation}
where $\omega^{(b)}_i = w^{\mathrm{bal}}_i\cdot m^{(b)}_i$ (or $w^{\mathrm{src}}_i m^{(b)}_i$), and $m^{(b)}_i$ is the
number of times row $i$ is drawn in the bootstrap sample of tree $b$ ($n$ draws with replacement). The balanced class
weights are computed once on the whole training fold, not per bootstrap sample. \textbf{Tree growth.} CART: at each node
draw $m=\max(1,\lfloor\sqrt p\rfloor)$ candidate features without replacement and choose the split that maximises the
weighted decrease of the Gini impurity
\begin{equation}
G(\mathcal N) = 1-\sum_{c\in\{0,1\}}\pi_c(\mathcal N)^2,\quad \pi_c(\mathcal N)=\frac{\sum_{i\in\mathcal N,y_i=c}\omega_i}{\sum_{i\in\mathcal N}\omega_i},\qquad
\Delta = \Omega_\mathcal{N} G(\mathcal N) - \Omega_{L}G(L) - \Omega_{R}G(R);\label{eq:gini}
\end{equation}
grow without depth limit until nodes are pure or a child would have fewer than \code{min\_samples\_leaf} rows.
With the two colours of H, $m=1$: every split uses one randomly chosen colour; with H plus three timing features and a
missing indicator ($p=6$), $m=2$. \textbf{Determinism.} \code{random\_state=42}. \textbf{Assumptions/limitations.}
Piecewise-constant, can follow non-monotone class regions (important for MAXI and for the timing$\times$colour
interaction), but with $<10$ BH sources a leaf can memorise one source; scores of trees are not calibrated (v3: A-RF had
good ranking but scores pushed below 0.5 -- a threshold problem).

\begin{algorithm}[htbp]
\caption{Random forest fit and score (one fold)}\label{alg:rf}
\begin{algorithmic}[1]
\Require $\Xm_\mathcal{T},\yv_\mathcal{T}$ (raw features), $\omega_i$ (balanced or source-equal), $B$, $m$, leaf size $L$
\For{$b=1,\dots,B$} \Comment{seeded}
  \State draw bootstrap counts $m^{(b)}\sim\mathrm{Multinomial}(n; 1/n,\dots,1/n)$; weights $\omega^{(b)}_i=\omega_i m^{(b)}_i$
  \State grow tree $b$: recursively split nodes on the best of $m$ random features by \cref{eq:gini}, leaves $\ge L$ rows
\EndFor
\State \Return $s_i$ by \cref{eq:rf} for $i\in\mathcal V$
\end{algorithmic}
\end{algorithm}

\begin{table}[htbp]
\centering\small
\caption[Random-forest hyper-parameters]{Random-forest hyper-parameters (\code{RF\_PARAMS}, \code{V4\_INNER\_TREES}, \code{V4\_OUTER\_TREES}).}\label{tab:hp-rf}
\begin{tabularx}{\textwidth}{@{}lp{3.6cm}lX@{}}
\toprule
Parameter & Value & Fixed / tuned & Grid / note \\
\midrule
\code{n\_estimators} & 500 & fixed & v4a nested: 200 trees in the inner loop, 500 for the outer refit \\
\code{max\_features} & \code{'sqrt'} & fixed & $m=\lfloor\sqrt p\rfloor$ \\
\code{min\_samples\_leaf} & 2 & fixed & v4a 1d/1e: $\{1,2,5,10\}$ \\
\code{class\_weight} & balanced (or \code{None} with $\wv^\mathrm{src}$) & fixed & -- \\
\code{bootstrap}, \code{criterion} & True, Gini & default & -- \\
\code{random\_state} & 42 & fixed & \code{n\_jobs}: 2 (v1--v3), 1 (\code{V4Model}) \\
\bottomrule
\end{tabularx}
\end{table}

% -----------------------------------------------------------------------------------------------------
\section{Extremely randomised trees}\label{sec:et}
% -----------------------------------------------------------------------------------------------------
\textbf{Role.} Exploratory alternative in v4a. \textbf{Difference from RF.} No bootstrap ($m^{(b)}_i=1$), and at each
node, for each of the $m$ candidate features one threshold is drawn uniformly between the feature's minimum and maximum
in the node; the best of these random splits by \cref{eq:gini} is used. Settings: 500 trees, \code{max\_features='sqrt'},
\code{min\_samples\_leaf=2}, balanced, \code{random\_state=42}; nested grid $\{1,2,5,10\}$ for the leaf size (200/500
trees inner/outer). Lower variance than RF at the price of more bias; no scaling needed.

\begin{algorithm}[htbp]
\caption{ExtraTrees fit and score}\label{alg:et}
\begin{algorithmic}[1]
\Require raw $\Xm_\mathcal T$, $\yv_\mathcal T$, weights $\omega_i$; $B=500$ (inner 200), $m=\lfloor\sqrt p\rfloor$, leaf size $L$
\For{$b=1,\dots,B$} \Comment{all rows, no bootstrap}
   \State at each node: draw $m$ features, one threshold $\sim U(\min,\max)$ each; keep the split with the largest $\Delta$ (\cref{eq:gini}); stop at purity or leaf size $L$
\EndFor
\State \Return $s_i=\frac1B\sum_b\hat p_b(\xv_i)$ (\cref{eq:rf} with $m^{(b)}_i=1$)
\end{algorithmic}
\end{algorithm}

\begin{table}[htbp]
\centering\small
\caption[ExtraTrees hyper-parameters]{ExtraTrees hyper-parameters (\code{V4\_FIXED['ExtraTrees']}, \code{V4\_GRID['ExtraTrees']}).}\label{tab:hp-et}
\begin{tabular}{@{}llll@{}}
\toprule
Parameter & Value & Fixed / tuned & Grid \\
\midrule
\code{n\_estimators} & 500 & fixed (inner 200) & -- \\
\code{max\_features} & \code{'sqrt'} & fixed & -- \\
\code{min\_samples\_leaf} & 2 & tuned in 1d/1e & $\{1,2,5,10\}$ \\
\code{class\_weight}, \code{random\_state} & balanced, 42 & fixed & -- \\
\bottomrule
\end{tabular}
\end{table}

% -----------------------------------------------------------------------------------------------------
\section{Histogram gradient boosting}\label{sec:histgb}
% -----------------------------------------------------------------------------------------------------
\textbf{Role.} Exploratory in v4a. \textbf{Model.} An additive model of $M$ regression trees on the log-odds scale,
$F_M(\xv) = F_0 + \eta\sum_{m=1}^{M} f_m(\xv)$, $s(\xv)=\sigma(F_M(\xv))$, with $F_0$ the weighted log-odds of the
training classes. Each tree $f_m$ is fitted to the weighted gradients and Hessians of the binary log loss at $F_{m-1}$,
$g_i = \omega_i\,(\sigma(F_{m-1}(\xv_i))-y_i)$, $h_i=\omega_i\,\sigma(1-\sigma)$; leaf values are
$-\sum_{i\in\mathrm{leaf}} g_i/(\sum_{i\in\mathrm{leaf}} h_i+\lambda)$ with $\lambda$ = \code{l2\_regularization} = 0; trees
are grown best-first on features pre-binned into at most 255 quantile bins, up to \code{max\_leaf\_nodes} leaves with at
least \code{min\_samples\_leaf} rows. Settings (\code{V4\_FIXED['HistGB']}): $\eta=0.1$, $M=200$, 15 leaves, 20 rows per
leaf, no early stopping, balanced, \code{random\_state=42}; nested grid $\eta\in\{0.05,0.1\}\times$ leaves $\in\{7,15\}$.
\textbf{Limitation.} With 30 training sources and 20 rows per leaf, a single source (15 rows) cannot form a leaf, which
partly protects against memorisation; with 2 colours the model is close to a smooth version of RF.

\begin{algorithm}[htbp]
\caption{HistGradientBoosting fit and score}\label{alg:histgb}
\begin{algorithmic}[1]
\Require raw $\Xm_\mathcal T$, $\yv_\mathcal T$, weights $\omega_i$; $\eta$, $M$, max leaves, min leaf rows
\State bin every feature into $\le255$ quantile bins (training rows); $F_0\gets\logit\bigl(\sum\omega_iy_i/\sum\omega_i\bigr)$
\For{$m=1,\dots,M$}
   \State $g_i\gets\omega_i(\sigma(F_{m-1}(\xv_i))-y_i)$, $h_i\gets\omega_i\sigma(1-\sigma)$
   \State grow a best-first tree on $(g,h)$; leaf value $-\sum g/(\sum h+\lambda)$; $F_m\gets F_{m-1}+\eta f_m$
\EndFor
\State \Return $s_i=\sigma(F_M(\xv_i))$
\end{algorithmic}
\end{algorithm}

\begin{table}[htbp]
\centering\small
\caption[HistGradientBoosting hyper-parameters]{HistGradientBoosting hyper-parameters (\code{V4\_FIXED['HistGB']}, \code{V4\_GRID['HistGB']}).}\label{tab:hp-histgb}
\begin{tabular}{@{}llll@{}}
\toprule
Parameter & Value & Fixed / tuned & Grid \\
\midrule
\code{learning\_rate} $\eta$ & 0.1 & tuned in 1d/1e & $\{0.05,0.1\}$ \\
\code{max\_leaf\_nodes} & 15 & tuned in 1d/1e & $\{7,15\}$ \\
\code{max\_iter} $M$ & 200 & fixed & -- \\
\code{min\_samples\_leaf} & 20 & fixed & -- \\
\code{l2\_regularization} $\lambda$ & 0 & fixed & -- \\
\code{early\_stopping} & False & fixed & -- \\
\code{class\_weight}, \code{random\_state} & balanced, 42 & fixed & -- \\
\bottomrule
\end{tabular}
\end{table}

% -----------------------------------------------------------------------------------------------------
\section{Support-vector machine with RBF kernel}\label{sec:svm}
% -----------------------------------------------------------------------------------------------------
\textbf{Role.} Exploratory in v4a (it was selected in 23 of 31 outer folds of the nested auto-selection) and descriptive
in v7c2; a Platt-scaled three-class version reproduces \citet{deBeurs2022} in v7c1. \textbf{Training} (libsvm dual, labels
$y'_i=2y_i-1$, standardised inputs $\vect z$):
\begin{equation}
\max_{\vect\alpha}\ \sum_i\alpha_i-\frac12\sum_{i,j}\alpha_i\alpha_jy'_iy'_jK(\vect z_i,\vect z_j)\quad
\text{s.t. } 0\le\alpha_i\le C\,c_{y_i},\ \sum_i\alpha_iy'_i=0,\qquad K(\vect z,\vect z')=e^{-\gamma\|\vect z-\vect z'\|^2},\label{eq:svm}
\end{equation}
with class weights $c_y = n/(2n_y)$. \textbf{Kernel width.} \code{gamma='scale'} sets $\gamma = 1/(p\,\Var(\vect Z))$,
where $\Var(\vect Z)$ is the variance of \emph{all} entries of the standardised training matrix ($\approx1$, so
$\gamma\approx1/p$); in the v4a grid $\gamma = \kappa/(p\Var(\vect Z))$ with factor $\kappa\in\{0.3,1,3\}$ (code
\code{gamma\_factor}). \textbf{Score.} The decision value $f(\vect z)=\sum_i\alpha_iy'_iK(\vect z_i,\vect z)+b$ is mapped by
the logistic function,
\begin{equation}
s(\xv) = \sigma\bigl(f(\vect z)\bigr),\label{eq:svmscore}
\end{equation}
so $s=0.5$ is exactly the SVM decision surface; Platt scaling is deliberately not used because it runs an internal random
cross-validation. \textbf{v7c2 variant.} $C\in\{0.1,1,10\}\times\gamma\in\{0.1,0.585,3\}$ (absolute $\gamma$),
standardised, balanced, same score; selected by inner grouped 5-fold source AUC (\cref{sec:nested}).
\textbf{v7c1 variant.} Three classes (BH, non-pulsing NS, pulsar), one-vs-one, $C=0.655298$, $\gamma=0.585144$ on
training-standardised features, unweighted, Platt probabilities (\code{probability=True, random\_state=42}), exactly as
the R/e1071 setting of \citet{deBeurs2022}.

\begin{algorithm}[htbp]
\caption{RBF-SVM fit and score (binary versions)}\label{alg:svm}
\begin{algorithmic}[1]
\Require $\Xm_\mathcal T$, $\yv_\mathcal T$; $C$; $\gamma$ or factor $\kappa$
\State standardise (\cref{eq:scaler}); if \code{'scale'}: $\gamma\gets\kappa/(p\,\Var(\vect Z_\mathcal T))$ ($\kappa=1$ for the fixed model)
\State solve the dual \cref{eq:svm} with box constraints $C\,c_{y_i}$ (libsvm SMO)
\State \Return $s_i=\sigma(f(\vect z_i))$ (\cref{eq:svmscore})
\end{algorithmic}
\end{algorithm}

\begin{table}[htbp]
\centering\small
\caption[SVM hyper-parameters]{SVM hyper-parameters (\code{V4\_FIXED['SVM\_RBF']}, \code{V4\_GRID}, \code{V7\_C2\_SVM\_GRID}, \code{V7\_DEBEURS\_SVM}).}\label{tab:hp-svm}
\begin{tabular}{@{}llll@{}}
\toprule
Use & $C$ & $\gamma$ & selection \\
\midrule
v4a fixed (1b) & 1 & \code{'scale'} & -- \\
v4a nested (1d/1e) & $\{0.1,1,10\}$ & $\{0.3,1,3\}\times$\code{'scale'} & inner LOSO, source AUC \\
v7c2 (MAXI) & $\{0.1,1,10\}$ & $\{0.1,0.585,3\}$ & inner grouped 5-fold, source AUC \\
v7c1 (reproduction) & 0.655298 & 0.585144 & fixed (from the paper), Platt, 3 classes \\
\bottomrule
\end{tabular}
\end{table}

% -----------------------------------------------------------------------------------------------------
\section{$k$-nearest-neighbour classifier}\label{sec:knn}
% -----------------------------------------------------------------------------------------------------
\textbf{Role.} Exploratory (v4a), descriptive (v7c2, v8d), reproduction (v7c1). \textbf{Model.} With $\mathcal N_k(\vect z)$
the $k$ training rows nearest to $\vect z$ in Euclidean distance on standardised features (uniform weights),
$\hat p(\xv)=\frac1k\sum_{i\in\mathcal N_k}y_i$. Because the training folds are imbalanced, the output is corrected to
equal class priors with $\pi = n_1/n$ of the training fold (\code{V4Model.score}, \code{C2Model.score}):
\begin{equation}
s(\xv) = \frac{\hat p/\pi}{\hat p/\pi + (1-\hat p)/(1-\pi)} .\label{eq:knnprior}
\end{equation}
\textbf{Hyper-parameters.} $k=15$ fixed (v4a 1b); grids $\{5,15,31\}$ (v4a nested) and $\{5,15,24,35\}$ (v7c2, v8d, inner
LOSO); $k$ is capped at the number of training rows. \textbf{v7c1.} Three classes, $k=24$, \emph{raw} (unstandardised)
features and plain vote fractions, as \citet{deBeurs2022}. \textbf{Limitation.} A test source's nearest neighbours can
all come from one training source of similar brightness or colour; this is exactly how the model can ``recognise'' a
source type rather than a compact-object type.

\begin{algorithm}[htbp]
\caption{kNN classifier (binary, prior-corrected)}\label{alg:knn}
\begin{algorithmic}[1]
\Require $\Xm_\mathcal T$, $\yv_\mathcal T$, $k$
\State standardise (\cref{eq:scaler}); $\pi\gets$ fraction of BH rows in $\mathcal T$; $k\gets\min(k,|\mathcal T|)$
\For{$i\in\mathcal V$} \State $\hat p\gets$ BH fraction among the $k$ nearest training rows; $s_i\gets$ \cref{eq:knnprior} \EndFor
\end{algorithmic}
\end{algorithm}

\begin{table}[htbp]
\centering\small
\caption[kNN hyper-parameters]{kNN hyper-parameters (\code{V4\_FIXED['kNN']}, \code{V4\_GRID['kNN']}, \code{V7\_C2\_KNN\_GRID}, \code{V7\_DEBEURS\_KNN\_K}).}\label{tab:hp-knn}
\begin{tabular}{@{}lllll@{}}
\toprule
Use & $k$ & scaling & output & selection \\
\midrule
v4a fixed & 15 & standardised & prior-corrected & -- \\
v4a nested & $\{5,15,31\}$ & standardised & prior-corrected & inner LOSO, source AUC \\
v7c2, v8d & $\{5,15,24,35\}$ & standardised & prior-corrected & inner LOSO, source AUC \\
v7c1 & 24 & raw & 3-class vote fractions & fixed (paper) \\
\bottomrule
\end{tabular}
\end{table}

% -----------------------------------------------------------------------------------------------------
\section{Quadratic discriminant analysis}\label{sec:qda}
% -----------------------------------------------------------------------------------------------------
\textbf{Role.} Exploratory (v4a). \textbf{Model.} Class-conditional Gaussians on standardised features with regularised
covariances and equal priors $\pi_0=\pi_1=\tfrac12$:
\begin{equation}
\hat\Sigma_c = (1-r)\,S_c + r\,\vect I,\qquad
s(\xv) = \frac{\pi_1\,\mathcal N(\vect z;\hat{\vect\mu}_1,\hat\Sigma_1)}{\sum_{c}\pi_c\,\mathcal N(\vect z;\hat{\vect\mu}_c,\hat\Sigma_c)},\label{eq:qda}
\end{equation}
where $S_c$ is the unbiased sample covariance of class $c$ (scikit-learn applies the shrinkage to the squared singular
values of the centred class data, which is equivalent). $r=0.1$ fixed; nested grid $\{0.01,0.1,0.5\}$ (\code{V4\_GRID['QDA']}).
QDA takes no sample weights, so class balance enters only through the priors. \textbf{Note.} Three QDA combinations were
the only v4a combinations whose source balanced-accuracy interval excluded 0 (\cref{ch:phase2}); with 76 combinations
this is within the expected number of chance exceedances.

\begin{algorithm}[htbp]
\caption{QDA fit and score}\label{alg:qda}
\begin{algorithmic}[1]
\Require $\Xm_\mathcal T$, $\yv_\mathcal T$, $r$
\State standardise; for $c\in\{0,1\}$: class mean $\hat{\vect\mu}_c$, unbiased covariance $S_c$, $\hat\Sigma_c\gets(1-r)S_c+r\vect I$
\State \Return $s_i$ by \cref{eq:qda} with $\pi_0=\pi_1=\tfrac12$
\end{algorithmic}
\end{algorithm}

\begin{table}[htbp]
\centering\small
\caption[QDA hyper-parameters]{QDA hyper-parameters (\code{V4\_FIXED['QDA']}, \code{V4\_GRID['QDA']}).}\label{tab:hp-qda}
\begin{tabular}{@{}llll@{}}
\toprule
Parameter & Value & Fixed / tuned & Grid \\
\midrule
\code{reg\_param} $r$ & 0.1 & tuned in 1d/1e & $\{0.01,0.1,0.5\}$ \\
\code{priors} & $[0.5,0.5]$ & fixed & -- \\
\bottomrule
\end{tabular}
\end{table}

% -----------------------------------------------------------------------------------------------------
\section{Multi-layer perceptron (seed ensemble)}\label{sec:mlp}
% -----------------------------------------------------------------------------------------------------
\textbf{Role.} Exploratory (v4a). \textbf{Model.} One hidden layer of 32 ReLU units (grid: $(16)$ or $(32,16)$), logistic
output, on standardised inputs. \textbf{Objective} (per minibatch $\mathcal M$; scikit-learn 1.8 with sample weights):
\begin{equation}
\mathcal L = \frac{\sum_{i\in\mathcal M}\omega_i\,\ell_i}{\sum_{i\in\mathcal M}\omega_i}+\frac{\alpha}{2}\,\frac{\sum_{W}\|W\|_F^2}{\sum_{i\in\mathcal M}\omega_i},\qquad \omega_i=w^\mathrm{bal}_i ,\label{eq:mlp}
\end{equation}
with $\ell_i$ the binary cross-entropy, $\alpha = 10^{-3}$ (grid $\{10^{-4},10^{-2}\}$); Adam (learning rate $10^{-3}$),
batch size $\min(200,n)$, at most 500 epochs. \textbf{Early stopping.} A stratified 20\% validation split (seeded with the
model seed) is held out; training stops when the weighted validation accuracy has not improved by $10^{-4}$ for 10 epochs,
and the weights with the best validation accuracy are restored. \textbf{Seed ensemble.} Because a small MLP is unstable,
the score is the mean of \code{predict\_proba} over seeds $\{0,\dots,4\}$ (outer fits) or $\{0,1,2\}$ (inner fits of the
nested procedure): $s(\xv)=\frac{1}{|\mathcal R|}\sum_{r\in\mathcal R}s_r(\xv)$. A planned 1D-CNN on representation B was
skipped as pre-declared because \texttt{torch} was not installed.

\begin{algorithm}[htbp]
\caption{MLP seed ensemble}\label{alg:mlp}
\begin{algorithmic}[1]
\Require $\Xm_\mathcal T$, $\yv_\mathcal T$; architecture, $\alpha$; seed set $\mathcal R$ (outer $\{0..4\}$, inner $\{0,1,2\}$)
\State standardise; $\omega_i\gets w^\mathrm{bal}_i$ (\cref{eq:balanced})
\For{$r\in\mathcal R$}
  \State stratified 80/20 split (seed $r$); Adam on \cref{eq:mlp}; after each epoch the weighted validation accuracy; stop after 10 epochs without improvement $>10^{-4}$ (max 500); restore the best weights
\EndFor
\State \Return $s_i=\frac{1}{|\mathcal R|}\sum_r s_r(\xv_i)$
\end{algorithmic}
\end{algorithm}

\begin{table}[htbp]
\centering\small
\caption[MLP hyper-parameters]{MLP hyper-parameters (\code{V4\_FIXED['MLP']}, \code{V4\_GRID['MLP']}, \code{V4\_MLP\_SEEDS\_*}).}\label{tab:hp-mlp}
\begin{tabularx}{\textwidth}{@{}p{4.6cm}Xll@{}}
\toprule
Parameter & Value & Fixed / tuned & Grid \\
\midrule
\code{hidden\_layer\_sizes} & $(32,)$ & tuned in 1d/1e & $\{(16,),(32,16)\}$ \\
\code{alpha} & $10^{-3}$ & tuned in 1d/1e & $\{10^{-4},10^{-2}\}$ \\
\code{early\_stopping}, \code{validation\_fraction} & True, 0.2 & fixed & -- \\
\code{max\_iter} & 500 & fixed & -- \\
activation, solver, learning rate, batch & ReLU, Adam, $10^{-3}$, $\min(200,n)$ & defaults & -- \\
seeds & outer $\{0,\dots,4\}$, inner $\{0,1,2\}$ & fixed & -- \\
\bottomrule
\end{tabularx}
\end{table}

% -----------------------------------------------------------------------------------------------------
\section{Colour-conditional timing likelihood-ratio classifier (CCTLR, v6c)}\label{sec:cctlr}
% -----------------------------------------------------------------------------------------------------
\textbf{Motivation.} In LR, a timing feature has one slope everywhere in the colour plane, although the timing information
is expected only where the colours overlap (hard states). CCTLR adds to the colour logit a generative log-likelihood ratio
of the timing features \emph{conditional on colour}, so that timing contributes only where the class-conditional timing
distributions differ, and missing timing is simply ``no evidence''. It was designed for v6 and has no tunable
hyper-parameter. \textbf{Model} (\code{v6lib.CCTLR}): with standardised colours $\vect h$ and timing vector
$\vect T=(T_1,T_2,T_3)$,
\begin{align}
\ell(\xv) &= \underbrace{\vect w^\top\vect h + b}_{\text{H-LR logit}} + \Lambda(\vect T\mid\vect h),\qquad s(\xv)=\sigma(\ell(\xv)),\label{eq:cctlr}\\
\Lambda(\vect T\mid\vect h) &= \sum_{k=1}^{3}\Bigl[\log t_4\!\bigl(T_k;\ \mu_{1k}(\vect h),\ \tau_{1k}\bigr)-\log t_4\!\bigl(T_k;\ \mu_{0k}(\vect h),\ \tau_{0k}\bigr)\Bigr], \nonumber\\
\mu_{ck}(\vect h) &= \beta_{ck0}+\beta_{ck1}h_1+\beta_{ck2}h_2, \nonumber
\end{align}
where $t_4(\cdot;\mu,\tau)$ is the Student-$t$ density with 4 degrees of freedom, location $\mu$ and scale $\tau$ (robust to
burst or dip outliers). For each class $c$ and timing feature $k$, $\vect\beta_{ck}$ is a weighted least-squares fit on
the training rows of class $c$ with finite timing, with source-equal weights within the class
$v_i = 1/(S_c\,n_{g_i})$ (\code{source\_equal\_within\_class}) and a ridge of $10^{-9}$; the scale is the weighted
residual standard deviation, floored at 0.01:
\begin{equation}
\hat{\vect\beta}_{ck}=\bigl(\vect Z_c^\top\vect V_c\vect Z_c+10^{-9}\vect I\bigr)^{-1}\vect Z_c^\top\vect V_c\vect T_{c,k},\qquad
\hat\tau_{ck}=\max\Bigl(0.01,\ \sqrt{\textstyle\sum_iv_ir_{ik}^2/\sum_iv_i}\Bigr),\label{eq:cctlrfit}
\end{equation}
with $\vect Z_c=[\vect 1,\vect h]$ of class $c$ and residuals $r_{ik}$. The LR part uses \code{LR\_PARAMS} (balanced, or the
source-equal weights of \cref{eq:sourceequal} in S3). If any timing value of a row is missing, $\Lambda=0$ and the score
equals the H-LR probability (unit test in \file{test\_v6lib.py}). \textbf{Limitation.} Adding a generative LR term to a
discriminative logit double counts nothing only if timing is conditionally independent of the label given colour under
the LR model -- an idealisation; the colour model of the class means is linear. In practice CCTLR did not beat H+T-LR
(\cref{ch:phase3}).

\begin{algorithm}[htbp]
\caption{CCTLR fit and score}\label{alg:cctlr}
\begin{algorithmic}[1]
\Require training colours $\vect H$, timing $\vect T$ (may contain NaN), labels $\yv$, sources $\gv$; test $(\vect H',\vect T')$
\State standardise $\vect H$ (\cref{eq:scaler}); fit H-LR (\cref{eq:lrobj}) on $\vect h$
\For{$c\in\{0,1\}$ and $k\in\{1,2,3\}$}
   \State rows: class $c$ with all timing finite; weights $v_i=1/(S_c n_{g_i})$; fit $\hat{\vect\beta}_{ck},\hat\tau_{ck}$ by \cref{eq:cctlrfit}
\EndFor
\For{each test row}
   \State $\Lambda\gets0$ if any $T'_k$ missing, else the sum in \cref{eq:cctlr}
   \State $s\gets\sigma(\vect w^\top\vect h'+b+\Lambda)$
\EndFor
\end{algorithmic}
\end{algorithm}

% -----------------------------------------------------------------------------------------------------
\section{$k$-nearest-neighbour regression for colour residualisation}\label{sec:knnreg}
% -----------------------------------------------------------------------------------------------------
\textbf{Role.} Inside every colour-conditional permutation test (v6d, v9b, v10b, v11b, v12a, v13b; \cref{sec:permutation}).
It removes the part of a timing feature that is predictable from colour \emph{without using labels}. For a feature
$F$ (e.g.\ $\log_{10}\nu_c$) and colours $\vect h$ standardised on the fit set (\code{v6lib.source\_residuals}):
\begin{equation}
\hat F_i = \frac1K\sum_{j\in\mathcal N_K(\vect h_i;\ \{j:\ g_j\ne g_i\})}F_j,\qquad r_i = F_i-\hat F_i,\qquad
\bar r_k=\frac{1}{|\{i:g_i=k\}|}\sum_{i:g_i=k} r_i ,\label{eq:knnreg}
\end{equation}
with $K=25$ (\code{V6\_KNN}; capped at the number of available rows) and uniform weights; the neighbours of an observation
are taken only from \emph{other} sources (leave-one-source-out), so a source cannot explain itself. Rows with missing $F$
are skipped. For the v6 external replication the regression is fitted on all S1 rows and applied to the external rows.
\textbf{Caveat} (noted in v9): when BHs cluster in one colour region, their neighbours are partly other BHs, so part of the
class difference is absorbed into $\hat F$; the test is then conservative.

\begin{algorithm}[htbp]
\caption{Label-free colour residualisation (\code{source\_residuals})}\label{alg:knnreg}
\begin{algorithmic}[1]
\Require colours $\vect H$, feature $F$ (may be NaN), sources $\gv$, $K=25$
\State standardise $\vect H$ on the fit set
\For{each source $k$} \State fit a kNN regression ($K$ neighbours) on rows with finite $F$ and $g\ne k$; predict rows of $k$; $r_i\gets F_i-\hat F_i$ \EndFor
\State \Return source means $\bar r_k$ (\cref{eq:knnreg})
\end{algorithmic}
\end{algorithm}

% -----------------------------------------------------------------------------------------------------
\section{Frozen characteristic-frequency rule (v11, v13)}\label{sec:nurule}
% -----------------------------------------------------------------------------------------------------
A deliberately simple, pre-frozen decision rule: an observation is called BH if it is hard-like and
$\nu_c<\nu^\ast$; a source is called BH if at least half of its new hard-like observations are called BH. The threshold
is the geometric mean of the class medians of $\nu_c$ in the previous data set,
\begin{equation}
\nu^\ast = \sqrt{\median_{\mathrm{BH}}\nu_c\cdot\median_{\mathrm{NS}}\nu_c},\label{eq:nustar}
\end{equation}
frozen before any new data were read: $\nu^\ast=\sqrt{0.5917\times2.8594}=1.3007$\,Hz from v10 (32 BH and 111 NS
hard-like observations; \file{logs/v11/nu\_star\_frozen.txt}) and $\nu^\ast=\sqrt{0.6294\times2.8362}=1.3361$\,Hz from the
v12 corrected and screened data (62 BH, 138 NS; \file{logs/v13/nu\_star\_frozen.txt}).

% -----------------------------------------------------------------------------------------------------
\section{Models that were not used}\label{sec:notused}
% -----------------------------------------------------------------------------------------------------
Principal component analysis was never used: every occurrence of ``PCA'' in the code refers to the RXTE Proportional
Counter Array. The pre-declared 1D convolutional network on representation B (v4) was skipped because \texttt{torch} was
not available; no replacement was introduced. No probability calibration (Platt, isotonic) was used except in the
de~Beurs reproduction (v7c1).
